\begin{table*}[t]
\caption{SWE-Bench Verified Results}
\centering
\renewcommand{\arraystretch}{1.3}
\resizebox{\textwidth}{!}{%
\begin{tabular}{c|c|c*{9}{|c}}
\toprule
\toprule
\multirow{2}{*}{\textbf{Scaffold}} & \multirow{2}{*}{\textbf{LLM}} & \multirow{2}{*}{\textbf{Training Algo.}} & \multicolumn{3}{|c}{\textbf{File-level}} & \multicolumn{3}{|c}{\textbf{Module-level}} & \multicolumn{3}{|c}{\textbf{Function-level}} \\
\cline{4-12}
& & & Recall (\%) & Precision (\%) & F1 (\%) & Recall (\%) & Precision (\%) & F1 (\%) & Recall (\%) & Precision (\%) & F1 (\%) \\
\midrule
\multicolumn{12}{c}{Closed-Source LLMs} \\
\midrule
RepoSearcher & Claude-3.7-Sonnet & - & 89.24 & 20.24 & 32.3 & - & - & - & 66.08 & 18.64 & 26.91 \\
\midrule
\multirow{3}{*}{RepoNavigator} & GPT-5-Chat & - & 58.17 & 61.87 & 58.88 & - & - & - & 30.42 & 34.56 & 31.17 \\
& Claude-3.7-Sonnet & - & 72.26 & 75.95 & 73.01 & - & - & - & 31.03 & 34.43 & 31.72 \\
& Claude-4.5-Sonnet & - & 80.68 & 81.92 & 79.94 & - & - & - & 43.97 & 45.76 & 43.62 \\
\midrule
\multicolumn{12}{c}{Open-Source LLMs} \\
\midrule
Agentless & \multirow{5}{*}{Qwen2.5-32B} & - & 78.93 & 25.6 & 35.38 & - & - & - & 40.97 & 24.07 & 27.33 \\
\cmidrule{1-1}\cmidrule{3-12}
CoSIL & & - & 83.5 & 19.34 & 30.77 & - & - & - & 55.38 & 14.85 & 22.11 \\
\cmidrule{1-1}\cmidrule{3-12}
LocAgent & & - & 79.39 & 34.18 & 44.18 & - & - & - & 46.79 & 16.29 & 21.48 \\
\cmidrule{1-1}\cmidrule{3-12}
OrcaLoca & & - & 59.57 & 59.51 & 58.11 & - & - & - & 39.14 & 25.59 & 28.72 \\
\midrule
\multirow{2}{*}{RepoSearcher} & Qwen2.5-7B & \multirow{2}{*}{Distillation+GRPO} & 83.11 & 18.8 & 30.09 & - & - & - & 62.38 & 17.68 & 25.57 \\
& Qwen2.5-32B &  & 88.59 & 20.24 & 32.25 & - & - & - & 68.55 & 19.36 & 28.03 \\
\midrule
\multirow{3}{*}{RepoNavigator} & Qwen2.5-7B & \multirow{3}{*}{GRPO} & 50.62 & 53.83 & 51.63 & - & - & - & 26.69 & 30.34 & 27.49 \\
& Qwen2.5-14B & & 61.6 & 58.97 & 58.9 & - & - & - & 31.02 & 30.08 & 29.23 \\
& Qwen2.5-32B & & 67.29 & 70.76 & 67.75 & - & - & - & 33.71 & 37.19 & 34.09 \\
\midrule
\multirow{5}{*}{OpenHands} & Qwen3-4B-Instruct & - & 53.34 & 49.69 & 49.73 & 20.15 & 19.86 & 19.32 & 13.74 & 14.17 & 13.27 \\
& Qwen3-14B & - & 71.2 & 36.49 & 43.13 & 33.04 & 20.4 & 22.86 & 23.58 & 14.51 & 16.08 \\
& Qwen3-32B & - & 70.41 & 58.63 & 60.88 & 40.27 & 35.57 & 35.95 & 28.69 & 25.94 & 25.77 \\
\cmidrule{2-12}
& \methodname-4B & \multirow{2}{*}{GSPO} & 67.74 & 71.53 & 68.52 & 44.97 & 49.7 & 45.97 & 35.72 & 40.71 & 36.78 \\
& \methodname-14B & & 68.69 & 71 & 68.57 & 50.88 & 53.71 & 50.88 & 40.27 & 43.74 & 40.32 \\
\bottomrule
\end{tabular}%
}
\label{tab:swe_bench_results}
\end{table*}
\section{Experimental Setup}
\label{sec:experiments}

In this section, we describe the detailed experimental setup used in our work.

\subsection{Training Setup}\label{sec:train_setup}

We prepare our training data by processing the SWE-smith~\citep{yang2025swesmithscalingdatasoftware} using the method described in \S\ref{sec:data_creation}. Our dataset consists of a total of 39K instances from 131 repositories after pre-processing. Note that none of these repositories are present in our evaluation benchmarks (\S\ref{sec:eval_benchmarks}) to avoid contamination.

We train models from the Qwen3 family~\citep{yang2025qwen3technicalreport}, particularly \texttt{Qwen/Qwen3-4B-Instruct-2507} and \texttt{Qwen/Qwen3-14B} (with thinking disabled). 

% \sonicomment{this can be shifted to appendix}
% \lscomment{I think this is essential detail given that we spent probably 2-3 weeks figuring out this particular issue}\sonicomment{ok}
We use a modified chat-template for the Qwen3-14B model that does not remove the \texttt{<think><\textbackslash think>} from earlier turns during tokenization~\footnote{https://huggingface.co/OpenPipe/Qwen3-14B-Instruct} rather than stripping them by default. This modification enables trajectory-level training, allowing us to concatenate the agent’s actions and the environment’s observations across turns into a single sequence, instead of training the model at an individual step level. We mask loss from the tokens which belong to the observations (i.e. tool responses) since these were not generated by the LLM. To avoid the retokenization drift~\footnote{https://blog.vllm.ai/2025/10/22/agent-lightning.html}, we obtain the prompt and response tokens directly from the vLLM engine. For all our training runs, we use the GSPO training algorithm as discussed in \S\ref{sec:train_algo}.

% \sonicomment{Given the difference in hyperparams across expts, shift the differences to appendix instead of highlighting them here.}
We train the 4B model for 200 steps, with a batch size of 8, sampling 8 rollouts per instance, and uses the reward design from \S\ref{sec:reward_design} but without the $r^{turn}$ reward, and allows a maximum of 6 turns during training. We train the 14B model for 300 steps with a batch size of 32, 4 rollouts per instance, uses the reward design from \S\ref{sec:reward_design}, and allows a maximum of 4 steps during training. For both training runs, we set the learning rate to $1e^{-6}$, \texttt{clip\_ratio\_low} to $3e^{-4}$, \texttt{clip\_ratio\_high} to $4e^{-4}$, temperature to 1.0. All training experiments were conducted on a single node consisting of 8xH100 GPUs.

\subsection{Evaluation Setup}\label{sec:eval_benchmarks}
We evaluate \methodname on three benchmarks: SWE-Bench Lite~\citep{jimenez2024swebenchlanguagemodelsresolve} consisting of 300 instances, SWE-Bench Verified~\footnote{https://openai.com/index/introducing-swe-bench-verified/} consisting of 500 instances, and the Python-subset of SWE-Bench Pro~\citep{deng2025swebenchproaiagents} consisting of 266 instances which consists of particularly difficult issues as compared to the other two benchmarks. We re-purpose these datasets and generate the ground truth using the approach from \S\ref{sec:data_creation}. During evaluation of Qwen3 models, we use the recommended generation hyper-parameters for instruct models/reasoning models with thinking disabled: temperature$=$0.7, topK$=$20, topP$=$0.8, and for reasoning models with thinking enabled: temperature$=$0.6, topP$=$0.95, topK$=$20.

\subsection{Evaluation Metrics}\label{sec:eval_metrics}
For each benchmark, we compute the F1 score for all three granularities (file, module, and function) and report their instance-wise averages. In addition to F1 score, we report the average precision and recall to compare with baselines that predict a fixed number of locations which may result in a high variance between the recall and precision.

\subsection{Baselines}\label{sec:baselines}
We compare \methodname against numerous competitive baselines: RepoNavigator~\citep{zhang2026toolenoughreinforcementlearning}, RepoSearcher~\citep{ma2025toolintegratedreinforcementlearningrepo}, OrcaLoca~\citep{yu2025orcalocallmagentframework}, CoSIL~\citep{liu2025cosil}, and Agentless~\citep{xia2024agentlessdemystifyingllmbasedsoftware} for different choices of LLMs. 

Note that all the above scaffolds except 
RepoNavigator predict a ranked list of top-$K$ files (where $K=5$ in general), as opposed to \methodname and RepoNavigator that let the model choose the locations dynamically. We hypothesize our design choice is better because higher precision is more crucial than higher recall for downstream issue resolution as \citet{cognition2025swegrep} find that polluting the context of the agent was more detrimental than leaving some context out, as the agent is typically only a few searches away to recover any remaining context. Furthermore, note that most of these approaches only report localization scores and file and function-level and not at the module-level.

Furthermore, we also evaluate the modified OpenHands~\citep{wang2025openhandssoftwareagentsdk} scaffold from \S\ref{sec:agent_scaffold} with the pre-trained versions of the models we train in our work: \texttt{Qwen/Qwen3-4B-Instruct-2507}, \texttt{Qwen/Qwen3-14B} (with OpenPipe chat template \S\ref{sec:train_setup}), and also with \texttt{Qwen/Qwen3-32B} with thinking enabled to compare our trained models against a larger-sized reasoning model.


